\title{2019 年高考数学 全国II卷-文科}
\date{}
\maketitle

\section*{一、选择题}
本题共 12 小题，每小题 5 分，共 60 分．在每小题给出的四个选项中，只有一项是符合题目要求的．

\begin{question}[points=5]
已知集合 $A=\{x\mid x>-1\}$，$B=\{x\mid x<2\}$，则 $A\cap B=$
\begin{choices}
  \item $(-1,+\infty)$
  \item $(-\infty,2)$
  \item $(-1,2)$
  \item $\varnothing$
\end{choices}
\begin{solution}
\textbf{答案：}$C$

\textbf{分析：}本题借助于数轴，根据交集的定义可得．

\textbf{详解：}由题知，$A\cap B=(-1,2)$，故选 $C$．
\end{solution}
\end{question}

\begin{question}[points=5]
设 $z=\mathrm{i}(2+\mathrm{i})$，则 $\overline{z}=$
\begin{choices}
  \item $1+2\mathrm{i}$
  \item $-1+2\mathrm{i}$
  \item $1-2\mathrm{i}$
  \item $-1-2\mathrm{i}$
\end{choices}
\begin{solution}
\textbf{答案：}$D$

\textbf{分析：}本题根据复数的乘法运算法则先求得 $z$，然后根据共轭复数的概念，写出 $\overline{z}$．

\textbf{详解：}$z=\mathrm{i}(2+\mathrm{i})=2\mathrm{i}+\mathrm{i}^2=-1+2\mathrm{i}$，所以 $\overline{z}=-1-2\mathrm{i}$，选 $D$．
\end{solution}
\end{question}

\begin{question}[points=5]
已知向量 $\vec{a}=(2,3)$，$\vec{b}=(3,2)$，则 $|\vec{a}-\vec{b}|=$
\begin{choices}
  \item $\sqrt{2}$
  \item $2$
  \item $5\sqrt{2}$
  \item $50$
\end{choices}
\begin{solution}
\textbf{答案：}$A$

\textbf{分析：}本题先计算 $\vec{a}-\vec{b}$，再根据模的概念求出 $|\vec{a}-\vec{b}|$．

\textbf{详解：}由已知，$\vec{a}-\vec{b}=(2,3)-(3,2)=(-1,1)$，所以 $|\vec{a}-\vec{b}|=\sqrt{(-1)^2+1^2}=\sqrt{2}$，故选 $A$．
\end{solution}
\end{question}

\begin{question}[points=5]
生物实验室有 $5$ 只兔子，其中只有 $3$ 只测量过某项指标，若从这 $5$ 只兔子中随机取出 $3$ 只，则恰有 $2$ 只测量过该指标的概率为
\begin{choices}
  \item $\dfrac{2}{3}$
  \item $\dfrac{3}{5}$
  \item $\dfrac{2}{5}$
  \item $\dfrac{1}{5}$
\end{choices}
\begin{solution}
\textbf{答案：}$B$

\textbf{分析：}本题首先用列举法写出所有基本事件，从中确定符合条件的基本事件数，应用古典概率的计算公式求解．

\textbf{详解：}设其中做过测试的 $3$ 只兔子为 $a,b,c$，剩余的 $2$ 只为 $A,B$，则从这 $5$ 只中任取 $3$ 只的所有取法有 $\{a,b,c\},\{a,b,A\},\{a,b,B\},\{a,c,A\},\{a,c,B\},\{a,A,B\},\{b,c,A\},\{b,c,B\},\{b,A,B\},\{c,A,B\}$ 共 $10$ 种．其中恰有 $2$ 只做过测试的取法有 $\{a,b,A\},\{a,b,B\},\{a,c,A\},\{a,c,B\},\{b,c,A\},\{b,c,B\}$ 共 $6$ 种，所以恰有 $2$ 只做过测试的概率为 $\dfrac{6}{10}=\dfrac{3}{5}$，选 $B$．
\end{solution}
\end{question}

\begin{question}[points=5]
在“一带一路”知识测验后，甲、乙、丙三人对成绩进行预测．

甲：我的成绩比乙高．

乙：丙的成绩比我和甲的都高．

丙：我的成绩比乙高．

成绩公布后，三人成绩互不相同且只有一个人预测正确，那么三人按成绩由高到低的次序为
\begin{choices}
  \item 甲、乙、丙
  \item 乙、甲、丙
  \item 丙、乙、甲
  \item 甲、丙、乙
\end{choices}
\begin{solution}
\textbf{答案：}$A$

\textbf{分析：}利用逐一验证的方法进行求解．

\textbf{详解：}若甲预测正确，则乙、丙预测错误，则甲比乙成绩高，丙比乙成绩低，故 $3$ 人成绩由高到低依次为甲，乙，丙；若乙预测正确，则丙预测也正确，不符合题意；若丙预测正确，则甲必预测错误，丙比乙的成绩高，乙比甲成绩高，即丙比甲、乙成绩都高，即乙预测正确，不符合题意，故选 $A$．
\end{solution}
\end{question}

\begin{question}[points=5]
设 $f(x)$ 为奇函数，且当 $x\geq 0$ 时，$f(x)=\mathrm{e}^x-1$，则当 $x<0$ 时，$f(x)=$
\begin{choices}
  \item $\mathrm{e}^{-x}-1$
  \item $\mathrm{e}^{-x}+1$
  \item $-\mathrm{e}^{-x}-1$
  \item $-\mathrm{e}^{-x}+1$
\end{choices}
\begin{solution}
\textbf{答案：}$D$

\textbf{分析：}先把 $x<0$，转化为 $-x>0$，代入可得 $f(-x)$，结合奇偶性可得 $f(x)$．

\textbf{详解：}$\because f(x)$ 是奇函数，当 $x<0$ 时，$-x>0$，$f(-x)=\mathrm{e}^{-x}-1=-f(x)$，得 $f(x)=-\mathrm{e}^{-x}+1$．故选 $D$．
\end{solution}
\end{question}

\begin{question}[points=5]
设 $\alpha,\beta$ 为两个平面，则 $\alpha\parallel\beta$ 的充要条件是
\begin{choices}
  \item $\alpha$ 内有无数条直线与 $\beta$ 平行
  \item $\alpha$ 内有两条相交直线与 $\beta$ 平行
  \item $\alpha,\beta$ 平行于同一条直线
  \item $\alpha,\beta$ 垂直于同一平面
\end{choices}
\begin{solution}
\textbf{答案：}$B$

\textbf{分析：}本题考查了空间两个平面的判定与性质及充要条件，渗透直观想象、逻辑推理素养，利用面面平行的判定定理与性质定理即可作出判断．

\textbf{详解：}由面面平行的判定定理知：$\alpha$ 内两条相交直线都与 $\beta$ 平行是 $\alpha/\!\/\beta$ 的充分条件，由面面平行性质定理知，若 $\alpha/\!\/\beta$，则 $\alpha$ 内任意一条直线都与 $\beta$ 平行，所以 $\alpha$ 内两条相交直线都与 $\beta$ 平行是 $\alpha/\!\/\beta$ 的必要条件，故选 $B$．
\end{solution}
\end{question}

\begin{question}[points=5]
若 $x_1=\dfrac{\pi}{4}$，$x_2=\dfrac{3\pi}{4}$ 是函数 $f(x)=\sin\omega x$ ($\omega>0$) 两个相邻的极值点，则 $\omega=$
\begin{choices}
  \item $2$
  \item $\dfrac{3}{2}$
  \item $1$
  \item $\dfrac{1}{2}$
\end{choices}
\begin{solution}
\textbf{答案：}$A$

\textbf{分析：}从极值点可得函数的周期，结合周期公式可得 $\omega$．

\textbf{详解：}由题意知，$f(x)=\sin\omega x$ 的周期 $T=\dfrac{2\pi}{\omega}=2\left(\dfrac{3\pi}{4}-\dfrac{\pi}{4}\right)=\pi$，得 $\omega=2$．故选 $A$．
\end{solution}
\end{question}

\begin{question}[points=5]
若抛物线 $y^2=2px$ ($p>0$) 的焦点是椭圆 $\dfrac{x^2}{3p}+\dfrac{y^2}{p}=1$ 的一个焦点，则 $p=$
\begin{choices}
  \item $2$
  \item $3$
  \item $4$
  \item $8$
\end{choices}
\begin{solution}
\textbf{答案：}$D$

\textbf{分析：}利用抛物线与椭圆有共同的焦点即可列出关于 $p$ 的方程，即可解出 $p$，或者利用检验排除的方法．

\textbf{详解：}因为抛物线 $y^2=2px(p>0)$ 的焦点 $\left(\dfrac{p}{2},0\right)$ 是椭圆 $\dfrac{x^2}{3p}+\dfrac{y^2}{p}=1$ 的一个焦点，所以 $3p-p=\left(\dfrac{p}{2}\right)^2$，解得 $p=8$，故选 $D$．
\end{solution}
\end{question}

\begin{question}[points=5]
曲线 $y=2\sin x+\cos x$ 在点 $(\pi,-1)$ 处的切线方程为
\begin{choices}
  \item $x-y-\pi-1=0$
  \item $2x-y-2\pi-1=0$
  \item $2x+y-2\pi+1=0$
  \item $x+y-\pi+1=0$
\end{choices}
\begin{solution}
\textbf{答案：}$C$

\textbf{分析：}先判定点 $(\pi,-1)$ 是否为切点，再利用导数的几何意义求解．

\textbf{详解：}当 $x=\pi$ 时，$y=2\sin\pi+\cos\pi=-1$，即点 $(\pi,-1)$ 在曲线 $y=2\sin x+\cos x$ 上．$y'=2\cos x-\sin x$，$y'|_{x=\pi}=2\cos\pi-\sin\pi=-2$，则 $y=2\sin x+\cos x$ 在点 $(\pi,-1)$ 处的切线方程为 $y-(-1)=-2(x-\pi)$，即 $2x+y-2\pi+1=0$．故选 $C$．
\end{solution}
\end{question}

\begin{question}[points=5]
已知 $a\in\left(0,\dfrac{\pi}{2}\right)$，$2\sin2\alpha=\cos2\alpha+1$，则 $\sin\alpha=$
\begin{choices}
  \item $\dfrac{1}{5}$
  \item $\dfrac{\sqrt{5}}{5}$
  \item $\dfrac{\sqrt{3}}{3}$
  \item $\dfrac{2\sqrt{5}}{5}$
\end{choices}
\begin{solution}
\textbf{答案：}$B$

\textbf{分析：}利用二倍角公式得到正余弦关系，利用角范围及正余弦平方和为 $1$ 关系得出答案．

\textbf{详解：}$\because 2\sin2\alpha=\cos2\alpha+1$，$\therefore 4\sin\alpha\cos\alpha=2\cos^2\alpha$．$\because \alpha\in\left(0,\dfrac{\pi}{2}\right)$，$\therefore \cos\alpha>0$，$\sin\alpha>0$，$\therefore 2\sin\alpha=\cos\alpha$，又 $\sin^2\alpha+\cos^2\alpha=1$，$\therefore 5\sin^2\alpha=1$，$\sin^2\alpha=\dfrac{1}{5}$，又 $\sin\alpha>0$，$\therefore \sin\alpha=\dfrac{\sqrt{5}}{5}$，故选 $B$．
\end{solution}
\end{question}

\begin{question}[points=5]
设 $F$ 为双曲线 $C:\dfrac{x^2}{a^2}-\dfrac{y^2}{b^2}=1$（$a>0$，$b>0$）的右焦点，$O$ 为坐标原点，以 $OF$ 为直径的圆与圆 $x^2+y^2=a^2$ 交于 $P$、$Q$ 两点．若 $|PQ|=|OF|$，则 $C$ 的离心率为
\begin{choices}
  \item $\sqrt{2}$
  \item $\sqrt{3}$
  \item $2$
  \item $\sqrt{5}$
\end{choices}
\begin{solution}
\textbf{答案：}$A$

\textbf{分析：}准确画图，由图形对称性得出 $P$ 点坐标，代入圆的方程得到 $c$ 与 $a$ 关系，可求双曲线的离心率．

\textbf{详解：}设 $PQ$ 与 $x$ 轴交于点 $A$，由对称性可知 $PQ\perp x$ 轴，又 $\because |PQ|=|OF|=c$，$\therefore |PA|=\dfrac{c}{2}$，$\therefore PA$ 为以 $OF$ 为直径的圆的半径，$\therefore A$ 为圆心，$|OA|=\dfrac{c}{2}$．$\therefore P\left(\dfrac{c}{2},\dfrac{c}{2}\right)$，又 $P$ 点在圆 $x^2+y^2=a^2$ 上，$\therefore \dfrac{c^2}{4}+\dfrac{c^2}{4}=a^2$，即 $\dfrac{c^2}{2}=a^2$，$\therefore e^2=\dfrac{c^2}{a^2}=2$，$\therefore e=\sqrt{2}$，故选 $A$．
\end{solution}
\end{question}

\section*{二、填空题}
本题共 4 小题，每小题 5 分，共 20 分．

\begin{question}[points=5]
若变量 $x,y$ 满足约束条件 $\begin{cases}2x+3y-6\ge 0,\\ x+y-3\le 0,\\ y-2\le 0,\end{cases}$ 则 $z=3x-y$ 的最大值是.
\fillin[{$9$}]
\begin{solution}
\textbf{答案：}$9$

\textbf{分析：}作出可行域，平移 $3x-y=0$ 找到目标函数取到最大值的点，求出点的坐标，代入目标函数可得．

\textbf{详解：}画出不等式组表示的可行域，阴影部分表示的三角形 $ABC$ 区域，根据直线 $3x-y-z=0$ 中的 $z$ 表示纵截距的相反数，当直线 $z=3x-y$ 过点 $C(3,0)$ 时，$z$ 取最大值为 $9$．
\end{solution}
\end{question}

\begin{question}[points=5]
我国高铁发展迅速，技术先进．经统计，在经停某站的高铁列车中，有 $10$ 个车次的正点率为 $0.97$，有 $20$ 个车次的正点率为 $0.98$，有 $10$ 个车次的正点率为 $0.99$，则经停该站高铁列车所有车次的平均正点率的估计值为.
\fillin[{$0.98$}]
\begin{solution}
\textbf{答案：}$0.98$

\textbf{分析：}本题考查通过统计数据进行概率的估计，采取估算法，利用概率思想解题．

\textbf{详解：}由题意得，经停该高铁站的列车正点数约为 $10\times0.97+20\times0.98+10\times0.99=39.2$，其中高铁个数为 $10+20+10=40$，所以该站所有高铁平均正点率约为 $\dfrac{39.2}{40}=0.98$．
\end{solution}
\end{question}

\begin{question}[points=5]
$\triangle ABC$ 的内角 $A,B,C$ 的对边分别为 $a,b,c$．已知 $b\sin A+a\cos B=0$，则 $B=$.
\fillin[{$\dfrac{3\pi}{4}$}]
\begin{solution}
\textbf{答案：}$\dfrac{3\pi}{4}$

\textbf{分析：}先根据正弦定理把边化为角，结合角的范围可得．

\textbf{详解：}由正弦定理，得 $\sin B\sin A+\sin A\cos B=0$．$\because A\in(0,\pi),B\in(0,\pi)$，$\therefore \sin A\neq 0$，得 $\sin B+\cos B=0$，即 $\tan B=-1$，$\therefore B=\dfrac{3\pi}{4}$．
\end{solution}
\end{question}

\begin{question}[points=5]
中国有悠久的金石文化，印信是金石文化的代表之一．印信的形状多为长方体、正方体或圆柱体，但南北朝时期的官员独孤信的印信形状是“半正多面体”（图 $1$）．半正多面体是由两种或两种以上的正多边形围成的多面体．半正多面体体现了数学的对称美．图 $2$ 是一个棱数为 $48$ 的半正多面体，它的所有顶点都在同一个正方体的表面上，且此正方体的棱长为 $1$．则该半正多面体共有个面，其棱长为．（本题第一空 $2$ 分，第二空 $3$ 分．）
\fillin[{$26$ 个面，$\sqrt{2}-1$}]
\begin{solution}
\textbf{答案：}$26$ 个面，$\sqrt{2}-1$

\textbf{分析：}第一问可按题目数出来，第二问需在正方体中简单还原出物体位置，利用对称性，平面几何解决．

\textbf{详解：}由图可知第一层与第三层各有 $9$ 个面，计 $18$ 个面，第二层共有 $8$ 个面，所以该半正多面体共有 $18+8=26$ 个面．如图，设该半正多面体的棱长为 $x$，则 $AB=BE=x$，延长 $BC$ 与 $FE$ 交于点 $G$，延长 $BC$ 交正方体棱于 $H$，由半正多面体对称性可知，$\triangle BGE$ 为等腰直角三角形，$\therefore BG=GE=CH=\dfrac{\sqrt{2}}{2}x$，$\therefore GH=2\times\dfrac{\sqrt{2}}{2}x+x=(\sqrt{2}+1)x=1$，$\therefore x=\dfrac{1}{\sqrt{2}+1}=\sqrt{2}-1$，即该半正多面体棱长为 $\sqrt{2}-1$．
\end{solution}
\end{question}

\section*{三、解答题}
共 70 分。解答应写出文字说明、证明过程或演算步骤。第 17～21 题为必考题，每个试题考生都必须作答．第 22、23 题为选考题，考生根据要求作答．

\begin{question}[points=12]
如图，长方体 $ABCD-A_1B_1C_1D_1$ 的底面 $ABCD$ 是正方形，点 $E$ 在棱 $AA_1$ 上，$BE\perp EC_1$．

（1）证明：$BE\perp$ 平面 $EB_1C_1$；

（2）若 $AE=A_1E$，$AB=3$，求四棱锥 $E-BB_1C_1C$ 的体积．

\begin{center}
  \includegraphics[width=.38\linewidth]{question-17-1.jpg}
\end{center}
\begin{solution}
\textbf{答案：}（1）见详解；（2）$18$

\textbf{分析：}（1）先由长方体得，$B_1C_1\perp$ 平面 $AA_1B_1B$，得到 $B_1C_1\perp BE$，再由 $BE\perp EC_1$，根据线面垂直的判定定理，即可证明结论成立；
（2）先设长方体侧棱长为 $2a$，根据题中条件求出 $a=3$；再取 $BB_1$ 中点 $F$，连结 $EF$，证明 $EF\perp$ 平面 $BB_1C_1C$，根据四棱锥的体积公式，即可求出结果．

\textbf{详解：}（1）因为在长方体 $ABCD-A_1B_1C_1D_1$ 中，$B_1C_1\perp$ 平面 $AA_1B_1B$；$BE\subset$ 平面 $AA_1B_1B$，所以 $B_1C_1\perp BE$，又 $BE\perp EC_1$，$B_1C_1\cap EC_1=C_1$，且 $EC_1\subset$ 平面 $EB_1C_1$，$B_1C_1\subset$ 平面 $EB_1C_1$，所以 $BE\perp$ 平面 $EB_1C_1$．
（2）设长方体侧棱长为 $2a$，则 $AE=A_1E=a$，由（1）可得 $EB_1\perp BE$；所以 $EB_1^2+BE^2=BB_1^2$，即 $2BE^2=BB_1^2$，又 $AB=3$，所以 $2AE^2+2AB^2=BB_1^2$，即 $2a^2+18=4a^2$，解得 $a=3$；取 $BB_1$ 中点 $F$，连结 $EF$，因为 $AE=A_1E$，则 $EF\parallel AB$；所以 $EF\perp$ 平面 $BB_1C_1C$，所以四棱锥 $E-BB_1C_1C$ 的体积为 $V_{E-BB_1C_1C}=\dfrac{1}{3}S_{\text{矩形}BB_1C_1C}\cdot EF=\dfrac{1}{3}\cdot BC\cdot BB_1\cdot EF=\dfrac{1}{3}\times3\times6\times3=18$．
\end{solution}
\end{question}

\begin{question}[points=12]
已知 $\{a_n\}$ 是各项均为正数的等比数列，$a_1=2$，$a_3=2a_2+16$．

（1）求 $\{a_n\}$ 的通项公式；

（2）设 $b_n=\log_2 a_n$，求数列 $\{b_n\}$ 的前 $n$ 项和．
\begin{solution}
\textbf{答案：}（1）$a_n=2^{2n-1}$；（2）$S_n=n^2$

\textbf{分析：}(1)本题首先可以根据数列 $\{a_n\}$ 是等比数列将 $a_3$ 转化为 $a_1q^2$，$a_2$ 转化为 $a_1q$，再然后将其带入 $a_3=2a_2+16$ 中，并根据数列 $\{a_n\}$ 是各项均为正数以及 $a_1=2$ 即可通过运算得出结果；
(2)本题可以通过数列 $\{a_n\}$ 的通项公式以及对数的相关性质计算出数列 $\{b_n\}$ 的通项公式，再通过数列 $\{b_n\}$ 的通项公式得知数列 $\{b_n\}$ 是等差数列，最后通过等差数列求和公式即可得出结果．

\textbf{详解：}(1)因为数列 $\{a_n\}$ 是各项均为正数的等比数列，$a_3=2a_2+16$，$a_1=2$，所以令数列 $\{a_n\}$ 的公比为 $q$，$a_3=a_1q^2=2q^2$，$a_2=a_1q=2q$，所以 $2q^2=4q+16$，解得 $q=-2$(舍去)或 $4$，所以数列 $\{a_n\}$ 是首项为 $2$、公比为 $4$ 的等比数列，$a_n=2\times4^{n-1}=2^{2n-1}$．
(2)因为 $b_n=\log_2 a_n$，所以 $b_n=2n-1$，$b_{n+1}=2n+1$，$b_{n+1}-b_n=2$，所以数列 $\{b_n\}$ 是首项为 $1$、公差为 $2$ 的等差数列，$S_n=\dfrac{1+(2n-1)}{2}\times n=n^2$．
\end{solution}
\end{question}

\begin{question}[points=12]
某行业主管部门为了解本行业中小企业的生产情况，随机调查了 $100$ 个企业，得到这些企业第一季度相对于前一年第一季度产值增长率 $y$ 的频数分布表．

\begin{center}

\begin{tabular}{c|ccccc}

$y$ 的分组 & $[-0.20,0)$ & $[0,0.20)$ & $[0.20,0.40)$ & $[0.40,0.60)$ & $[0.60,0.80)$ \\ \hline

企业数 & $2$ & $24$ & $53$ & $14$ & $7$

\end{tabular}

\end{center}

（1）分别估计这类企业中产值增长率不低于 $40\%$ 的企业比例、产值负增长的企业比例；



（2）求这类企业产值增长率的平均数与标准差的估计值（同一组中的数据用该组区间的中点值为代表）．（精确到 $0.01$）



附：$\sqrt{74}\approx8.602$．
\begin{solution}
\textbf{答案：}（1）增长率不低于 $40\%$ 的企业比例为 $21\%$，产值负增长的企业比例为 $2\%$；（2）平均数 $0.30$，标准差 $0.17$

\textbf{分析：}(1)本题首先可以通过题意确定 $100$ 个企业中增长率不低于 $40\%$ 的企业以及产值负增长的企业的个数，然后通过增长率不低于 $40\%$ 的企业以及产值负增长的企业的个数除随机调查的企业总数即可得出结果；
(2)可通过平均值以及标准差的计算公式得出结果．

\textbf{详解：}(1)由题意可知，随机调查的 $100$ 个企业中增长率不低于 $40\%$ 的企业有 $14+7=21$ 个，产值负增长的企业有 $2$ 个，所以增长率不低于 $40\%$ 的企业比例为 $\dfrac{21}{100}=0.21$，产值负增长的企业比例为 $\dfrac{2}{100}=0.02$．
(2)由题意可知，平均值 $\overline{y}=\dfrac{2\times(-0.10)+24\times0.10+53\times0.30+14\times0.50+7\times0.70}{100}=0.30$，标准差的平方：$s^2=\dfrac{1}{100}[2\times(-0.10-0.30)^2+24\times(0.10-0.30)^2+53\times(0.30-0.30)^2+14\times(0.50-0.30)^2+7\times(0.70-0.30)^2]=\dfrac{1}{100}[0.32+0.96+0+0.56+1.12]=0.0296$，所以标准差 $s=\sqrt{0.0296}=0.02\times\sqrt{74}\approx0.17$．
\end{solution}
\end{question}

\begin{question}[points=12]
已知 $F_1,F_2$ 是椭圆 $C:\dfrac{x^2}{a^2}+\dfrac{y^2}{b^2}=1(a>b>0)$ 的两个焦点，$P$ 为 $C$ 上一点，$O$ 为坐标原点．

（1）若 $\triangle POF_2$ 为等边三角形，求 $C$ 的离心率；

（2）如果存在点 $P$，使得 $PF_1\perp PF_2$，且 $\triangle F_1PF_2$ 的面积等于 $16$，求 $b$ 的值和 $a$ 的取值范围．
\begin{solution}
\textbf{答案：}（1）$e=\sqrt{3}-1$；（2）$b=4$，$a$ 的取值范围为 $[4\sqrt{2},+\infty)$

\textbf{分析：}（1）先连结 $PF_1$，由 $\triangle POF_2$ 为等边三角形，得到 $\angle F_1PF_2=90^\circ$，$PF_2=c$，$PF_1=\sqrt{3}c$；再由椭圆定义，即可求出结果；
（2）先由题意得到，满足条件的点 $P(x,y)$ 存在，当且仅当 $\dfrac{1}{2}|y|\cdot2c=16$，$\dfrac{y}{x+c}\cdot\dfrac{y}{x-c}=-1$，$\dfrac{x^2}{a^2}+\dfrac{y^2}{b^2}=1$，根据三个式子联立，结合题中条件，即可求出结果．

\textbf{详解：}（1）连结 $PF_1$，由 $\triangle POF_2$ 为等边三角形可知：在 $\triangle F_1PF_2$ 中，$\angle F_1PF_2=90^\circ$，$PF_2=c$，$PF_1=\sqrt{3}c$，于是 $2a=PF_1+PF_2=c+\sqrt{3}c$，故椭圆 $C$ 的离心率为 $e=\dfrac{c}{a}=\dfrac{1}{1+\sqrt{3}}=\sqrt{3}-1$．
（2）由题意可知，满足条件的点 $P(x,y)$ 存在，当且仅当 $\dfrac{1}{2}|y|\cdot2c=16$，$\dfrac{y}{x+c}\cdot\dfrac{y}{x-c}=-1$，$\dfrac{x^2}{a^2}+\dfrac{y^2}{b^2}=1$，即 $c|y|=16$ ①，$x^2+y^2=c^2$ ②，$\dfrac{x^2}{a^2}+\dfrac{y^2}{b^2}=1$ ③，由②③以及 $a^2=b^2+c^2$ 得 $y^2=\dfrac{b^4}{c^2}$，又由①知 $y^2=\dfrac{16^2}{c^2}$，故 $b=4$；由②③得 $x^2=\dfrac{a^2}{c^2}(c^2-b^2)$，所以 $c^2\ge b^2$，从而 $a^2=b^2+c^2\ge 2b^2=32$，故 $a\ge 4\sqrt{2}$；当 $b=4$，$a\ge 4\sqrt{2}$ 时，存在满足条件的点 $P$．所以 $b=4$，$a$ 的取值范围为 $[4\sqrt{2},+\infty)$．
\end{solution}
\end{question}

\begin{question}[points=12]
已知函数 $f(x)=(x-1)\ln x-x-1$．证明：

（1）$f(x)$ 存在唯一的极值点；

（2）$f(x)=0$ 有且仅有两个实根，且两个实根互为倒数．
\begin{solution}
\textbf{答案：}（1）见详解；（2）见详解

\textbf{分析：}（1）先对函数 $f(x)$ 求导，根据导函数的单调性，得到存在唯一 $x_0$，使得 $f'(x_0)=0$，进而判断函数 $f(x)$ 的单调性，即可确定其极值点个数，证明出结论成立；
（2）先由（1）的结果，得到 $f(x_0)<f(1)=-2<0$，$f(e^2)=e^2-3>0$，得到 $f(x)=0$ 在 $(x_0,+\infty)$ 内存在唯一实根，记作 $x=\alpha$，再求出 $f\left(\dfrac{1}{\alpha}\right)=0$，即可结合题意，说明结论成立．

\textbf{详解：}（1）由题意可得，$f(x)$ 的定义域为 $(0,+\infty)$，由 $f(x)=(x-1)\ln x-x-1$，得 $f'(x)=\ln x+\dfrac{x-1}{x}-1=\ln x-\dfrac{1}{x}$，显然 $f'(x)=\ln x-\dfrac{1}{x}$ 单调递增；又 $f'(1)=-1<0$，$f'(2)=\ln2-\dfrac{1}{2}=\dfrac{\ln4-1}{2}>0$，故存在唯一 $x_0$，使得 $f'(x_0)=0$；又当 $x>x_0$ 时，$f'(x)>0$，函数 $f(x)$ 单调递增；当 $0<x<x_0$ 时，$f'(x)<0$，函数 $f(x)$ 单调递减；因此，$f(x)$ 存在唯一的极值点．
（2）由（1）知，$f(x_0)<f(1)=-2$，又 $f(e^2)=e^2-3>0$，所以 $f(x)=0$ 在 $(x_0,+\infty)$ 内存在唯一实根，记作 $x=\alpha$．由 $1<x_0<\alpha$ 得 $\dfrac{1}{\alpha}<1<x_0$，又 $f\left(\dfrac{1}{\alpha}\right)=\left(\dfrac{1}{\alpha}-1\right)\ln\dfrac{1}{\alpha}-\dfrac{1}{\alpha}-1=\dfrac{f(\alpha)}{\alpha}=0$，故 $\dfrac{1}{\alpha}$ 是方程 $f(x)=0$ 在 $(0,x_0)$ 内的唯一实根；综上，$f(x)=0$ 有且仅有两个实根，且两个实根互为倒数．
\end{solution}
\end{question}

\section*{（二）选考题}
共 10 分．请考生在第 22、23 题中任选一题作答。如果多做，则按所做的第一题计分．

\begin{question}[points=10]
[选修 4-4：坐标系与参数方程]

在极坐标系中，$O$ 为极点，点 $M(\rho_0,\theta_0)(\rho_0>0)$ 在曲线 $C:\rho=4\sin\theta$ 上，直线 $l$ 过点 $A(4,0)$ 且与 $OM$ 垂直，垂足为 $P$．

（1）当 $\theta_0=\dfrac{\pi}{3}$ 时，求 $\rho_0$ 及 $l$ 的极坐标方程；

（2）当 $M$ 在 $C$ 上运动且 $P$ 在线段 $OM$ 上时，求 $P$ 点轨迹的极坐标方程．
\begin{solution}
\textbf{答案：}（1）$\rho_0=2\sqrt{3}$，$l$ 的极坐标方程为 $\rho\cos\left(\theta-\dfrac{\pi}{3}\right)=2$ 或 $\rho\sin\left(\theta+\dfrac{\pi}{6}\right)=2$；（2）$\rho=4\cos\theta\left(\dfrac{\pi}{4}\le\theta\le\dfrac{\pi}{2}\right)$

\textbf{分析：}（1）先由题意，将 $\theta_0=\dfrac{\pi}{3}$ 代入 $\rho=4\sin\theta$ 即可求出 $\rho_0$；根据题意求出直线 $l$ 的直角坐标方程，再化为极坐标方程即可；
（2）先由题意得到 $P$ 点轨迹的直角坐标方程，再化为极坐标方程即可，要注意变量的取值范围．

\textbf{详解：}（1）因为点 $M(\rho_0,\theta_0)(\rho_0>0)$ 在曲线 $C:\rho=4\sin\theta$ 上，所以 $\rho_0=4\sin\theta_0=4\sin\dfrac{\pi}{3}=2\sqrt{3}$；即 $M\left(2\sqrt{3},\dfrac{\pi}{3}\right)$，所以 $k_{OM}=\tan\dfrac{\pi}{3}=\sqrt{3}$，因为直线 $l$ 过点 $A(4,0)$ 且与 $OM$ 垂直，所以直线 $l$ 的直角坐标方程为 $y=-\dfrac{\sqrt{3}}{3}(x-4)$，即 $x+\sqrt{3}y-4=0$；因此，其极坐标方程为 $\rho\cos\theta+\sqrt{3}\rho\sin\theta=4$，即 $\rho\cos\left(\theta-\dfrac{\pi}{3}\right)=2$ 或 $\rho\sin\left(\theta+\dfrac{\pi}{6}\right)=2$．
（2）设 $P(x,y)$，则 $k_{OP}=\dfrac{y}{x}$，$k_{AP}=\dfrac{y}{x-4}$，由题意，$OP\perp AP$，所以 $k_{OP}k_{AP}=-1$，故 $\dfrac{y^2}{x^2-4x}=-1$，整理得 $x^2+y^2-4x=0$，因为 $P$ 在线段 $OM$ 上，$M$ 在 $C$ 上运动，所以 $0\le x\le2$，$2\le y\le4$，所以 $P$ 点轨迹的极坐标方程为 $\rho^2-4\rho\cos\theta=0$，即 $\rho=4\cos\theta$，$\dfrac{\pi}{4}\le\theta\le\dfrac{\pi}{2}$．
\end{solution}
\end{question}

\begin{question}[points=10]
[选修 4-5：不等式选讲]

已知 $f(x)=|x-a|x+|x-2|(x-a)$．

（1）当 $a=1$ 时，求不等式 $f(x)<0$ 的解集；

（2）若 $x\in(-\infty,1)$ 时，$f(x)<0$，求 $a$ 的取值范围．
\begin{solution}
\textbf{答案：}（1）$(-\infty,1)$；（2）$[1,+\infty)$

\textbf{分析：}（1）根据 $a=1$，将原不等式化为 $|x-1|x+|x-2|(x-1)<0$，分别讨论 $x<1$，$1\le x<2$，$x\ge2$ 三种情况，即可求出结果；
（2）分别讨论 $a\ge1$ 和 $a<1$ 两种情况，即可得出结果．

\textbf{详解：}（1）当 $a=1$ 时，原不等式可化为 $|x-1|x+|x-2|(x-1)<0$；当 $x<1$ 时，原不等式可化为 $(1-x)x+(2-x)(x-1)<0$，即 $(x-1)^2>0$，显然成立，此时解集为 $(-\infty,1)$；当 $1\le x<2$ 时，原不等式可化为 $(x-1)x+(2-x)(x-1)<0$，解得 $x<1$，此时解集为空集；当 $x\ge2$ 时，原不等式可化为 $(x-1)x+(x-2)(x-1)<0$，即 $(x-1)^2<0$，显然不成立；此时解集为空集；综上，原不等式的解集为 $(-\infty,1)$．
（2）当 $a\ge1$ 时，因为 $x\in(-\infty,1)$，所以由 $f(x)<0$ 可得 $(a-x)x+(2-x)(x-a)<0$，即 $(x-a)(x-1)>0$，显然恒成立；所以 $a\ge1$ 满足题意；当 $a<1$ 时，$f(x)=\begin{cases}2(x-a), & a\le x<1 \\ 2(x-a)(1-x), & x<a\end{cases}$，因为 $a\le x<1$ 时，$f(x)<0$ 显然不能成立，所以 $a<1$ 不满足题意；综上，$a$ 的取值范围是 $[1,+\infty)$．
\end{solution}
\end{question}

